\chapter*{Introduction}
\addcontentsline{toc}{chapter}{Introduction}

The LINKU camera payload measures a thin tether strung between two satellites.
It returns metric scalars about that tether, its arc length, straightness,
depth range, diameter and curvature, and a sparse point cloud as a by-product.
Every stage of the processing chain has been built and exercised, and the
optical and photometric requirements have been derived from primary sources in
the companion deliverables. What has not happened is the thing all of it exists
for: the instrument has never been pointed at a real tether, in the dark, on
the support it will actually fly on, with the tether and the rig both in
motion.

That is this experiment. In the project work plan \cite{linku_workplan} it is
the package \emph{Experimento de caja negra, hilo en movimiento}, inside the
phase-2 deliverable \emph{Integracion y experimento en la estructura LINKU}. It
is run in a blacked-out television studio of the PUCP Faculty of
Communications, referred to throughout this document as \textbf{TV2}, a space
borrowed for the purpose because it offers what the project needs and does not
own: a room that can be made genuinely dark, with professional lighting
fixtures and the rigging to position them.

\section*{Why this document exists}

Three reasons, in order of importance.

\textbf{First, the acceptance criteria have to be written down before the data
exists.} This is not a point of style. A review of the project's own
deliverables shows the same gap in three of them independently: the
rolling-shutter experiment and the Blender validation each carry a pending item
instructing a future author to define the quantitative agreement threshold
before running, rather than after seeing the results; the calibration method
leaves three of its four acceptance criteria numerically unset; and the camera
system report \cite{linku_d5_camera_report} contains no reconstruction accuracy
figure at all, reporting its two reconstruction scenarios qualitatively. The
common root is that no mission accuracy budget for tether reconstruction has
been written anywhere in the project. This document cannot supply that budget,
but it can refuse to inherit the habit. Chapter \ref{chap:criteria} registers
the criteria for this experiment in advance.

\textbf{Second, the parameters are scattered and some of them are
destructive to get wrong.} The illumination requirements live in one
deliverable, the optical discipline in another, the calibration preconditions
in a third, and the angular-rate limits in a fourth. On the day, in a borrowed
room, with a time slot, nobody is going to open four documents. Chapter
\ref{chap:parameters} collects all of them in one place, each with its
derivation and its source, and marks which are physics and which are project
decisions open to revision.

\textbf{Third, the session has to be replicable in Blender afterwards.} The
Blender validation deliverable \cite{linku_blender_validation} compares a real
photograph against a render of the same scene, and it currently cannot proceed
because the laboratory-condition light values in its scene are placeholders and
the physical pose of the tether in the chosen reference pair was never written
down. Chapter \ref{chap:record} defines the record so that this experiment does
not repeat that omission.

\section*{What this document is not}

It is not a specification of the illumination; that is derived from primary
sources in the dark-room deliverable \cite{linku_darkroom_illumination} and is
only summarised here. It is not a calibration procedure; that is specified and
implemented in the calibration deliverable \cite{linku_calibration_method} and
is invoked here as a precondition. It is not an analysis of rotation effects
with a known angular rate; that requires a motorised stage and is a separate
experiment \cite{linku_rollingshutter}. In TV2 both the tether and the rig are
moved \emph{by hand}, erratically and without instrumentation of the motion,
which is a deliberate choice discussed in Section \ref{sec:why-by-hand}.

\section*{Status of this revision}

Chapters \ref{chap:objectives} through \ref{chap:analysis} are complete and do
not depend on results. Chapter \ref{chap:results} is empty: at the time of
writing the repository contains no capture from this session, the most recent
non-simulated stereo pairs dating from February 2026, and the record sheet of
the pre-session parameter artifact has not been filled in. Chapter
\ref{chap:results} states precisely what goes in it.
